\documentclass[UTF8,10pt,a4paper]{ctexart}
\usepackage[margin=1.65cm,headheight=14pt]{geometry}
\usepackage{amsmath,booktabs,array,graphicx,xcolor,hyperref,xurl,fancyhdr}
\setmainfont{Arial}\setsansfont{Arial}\setmonofont{Menlo}
\setCJKmainfont{Songti SC}\setCJKsansfont{Heiti SC}
\definecolor{ink}{HTML}{163C50}\definecolor{accent}{HTML}{007C91}
\hypersetup{colorlinks=true,linkcolor=ink,urlcolor=accent,pdftitle={S87有限末端对照：实际结果与零基础讲解}}
\setlength{\parskip}{.55em}\setlength{\parindent}{1.5em}\linespread{1.07}
\setlength{\emergencystretch}{2em}\renewcommand{\arraystretch}{1.17}
\pagestyle{fancy}\fancyhf{}\lhead{\small S87：实际结果与零基础讲解}\rhead{\small 2026-09-11}\cfoot{\thepage}
\ctexset{section={format=\Large\bfseries\sffamily\color{ink},beforeskip=0pt,afterskip=.7em}}
\newcommand{\boxnote}[1]{\begin{center}\fcolorbox{accent}{accent!5}{\parbox{.92\linewidth}{#1}}\end{center}}
\begin{document}
\section*{1. 这轮回答了什么？}
\boxnote{\textbf{已有普通末端策略达到或低于旧多步引导的本例MSE；“取得该分数必须多步引导”的说法被反例否定。} 本轮是已有真实结果上的有限派生实验，没有新增完整生成链，也没有验证新方法。}
本轮生成完成时间：2026-09-11 07:38:47（北京时间）。随后单次固定评分和团队内不同作者复算均通过。旧S86的16条记录保留为历史引用；本轮只有6策略、24条新目标记录，不能说40个新独立样本。
本例达到反例标准的Gterminal 0.75只在四帧平均上更低；目标20、21、23较低，目标22反而由旧Gguide的0.0682643877变为0.07483243左右。全部逐目标数据仍保留，不能说四张图逐一改善。
\begin{center}\small\begin{tabular}{lrrr}\toprule
本轮固定策略 & 四目标总SSE & 全图MSE & 相对旧Gguide差\\\midrule
图像混合 0.5 & 16,836,148,626 & 0.0650333577 & 0.0126111351\\
末步融合 0.5 & 17,489,680,326 & 0.0675577688 & 0.0151355463\\
图像混合 0.75 & 13,814,266,216 & 0.0533606667 & 0.0009384442\\
末步融合 0.75 & 13,246,509,800 & 0.0511675817 & -0.0012546409\\
图像混合 1.0 & 14,513,028,466 & 0.0560597909 & 0.0036375683\\
末步融合 1.0 & 13,599,748,865 & 0.0525320459 & 0.0001098234\\
\bottomrule\end{tabular}\end{center}
\begin{center}\small\begin{tabular}{lrl}\toprule
S86历史参照 & 全图MSE & 说明\\\midrule
G0 & 0.1311666678 & 原始完整链\\
Gpaste & 0.0909680125 & 旧末端强度0.25\\
Gterminal & 0.0981916024 & 旧末端强度0.25\\
Gguide & 0.0524222225 & 50步引导，强度0.25\\
\bottomrule\end{tabular}\end{center}
主指标是每个对应颜色标量的平方差，除以颜色标量总数与$255^2$。每个策略使用相同四帧、同一强度；全图均值越小，只说明这个指标越低。

表中差为“新策略减旧Gguide”，负号表示新策略的误差更低。最终判断用未四舍五入的整数SSE：新策略是否不大于13,571,317,266。不能靠表中小数显示的近似平局作决定。

这是已见静态场景上的事后有限对照。先固定三个强度再运行，能够约束这次比较，但不能把之前看过的场景重新变成未见验证集。
\clearpage
\section*{2. 它处在proposal的哪一步？}

proposal研究的是：当系统见过一个三维环境，再换视角、过一段时间或再次回到某处时，生成内容能否继续保持一致。位置、形状、相机视角和动态状态都与此有关。最终需要处理长期和动态情况；这四张静态目标只是把一个失败拆开检查。

\textbf{第一步：建立能真实运行的普通基线。}之前已经运行原VMem相关采样与权重，并明确使用ft-mse VAE变体。VAE是把压缩表示还原成图像的组件；其原SD2.1身份仍未确认，不能说完整原系统的每一部分都精确复现。

\textbf{第二步：把“记住了”和“用对了”分开。}选到历史照片，不等于在请求的相机位置正确使用它。之前先检查相机与对应、历史预测几何和投影，再把同一个固定历史投影交给生成端。

\textbf{第三步：用普通对照排除容易误判的解释。}S86比较原生成、最后在RGB图像混合、最后在内部表示融合，以及50步采样引导。多步引导的MSE更低，但图上仍有重影；同时50次与1次的累计影响不同。

\textbf{当前S87：检查“末端对照只是太弱”这一解释。}原末端强度只有0.25。本轮固定0.5、0.75、1，检查普通末端办法能否达到那次多步引导的像素分数。这是在测试机制结论是否站得住，不是把更强的混合重新命名为创新。

\textbf{后续真正的方法阶段仍缺什么？}首先需要普通强方法确实解决不了、且可以重复观察的失败；其次要有不同于近邻工作的机制；最后用未用于选规则的场景、清楚的相机或物体参考、消融和代价比较检验。当前没有选定或验证新方法，也没有长期动态与跨场景确认。

\boxnote{阶段编号、运行次数、日志页数和模型花费不等于proposal完成百分比，更不能直接换算学生工时或PhD水平。当前可以准确汇报的是完成了哪些真实对照、排除了哪些说法、还缺哪类证据。}

\clearpage
\section*{3. 两类处理：同一个强度为什么不等价？}

\textbf{共同输入。}四张历史图编号19、18、13、12，四目标20、21、22、23。完整8槽先历史后目标。历史投影与两种mask均沿用S86保存量，没有重新估计几何，也没有从评分参考中补颜色或深度。

\textbf{图像混合Gpaste。}先把G0原始浮点输出按原规则转成$[0,1]$并截断范围，再令$w=\lambda M$，$M$是576像素图像上的支持区域：
\[I'=(1-w)I_0+wW_{RGB}\quad(w>0),\qquad I'=I_0\quad(w=0).\]
使用浮点图混合，最后才转成uint8。这里是普通RGB层面的合成，没有运行新神经网络；孔洞保留G0。强度1只在有支持的位置直接用历史投影颜色，投影本身未被保证正确。

\textbf{末步融合Gterminal。}从G0实际保存的最后一步状态出发，对内部压缩表示的clean预测$d$作融合：
\[d'=(1-\lambda m)d+\lambda mW.\]
其中$m$是每个$8\times8$图像块的平均支持比例，不是预测置信度。历史槽和零支持位置保持原值。随后沿原FP32顺序计算：
\[v=(x_{\mathrm{tilde}}-d')/\widehat\sigma,\qquad
z'=x_{\mathrm{tilde}}+(\sigma_{\mathrm{next}}-\widehat\sigma)v.\]
保留原$\widehat\sigma$中的小常数与计算顺序，不能在代码里用“最后应等于$d'$”跳过实际Euler算术。每个强度都完整解码8槽，再取后4张目标。

\textbf{为何要保留这些细节？}当$\lambda=1,m=0.5$时，内部表示的实际混合权重仍是0.5。VAE解码又是非线性并会传播局部变化，因此相同$\lambda$不代表相同最终RGB影响，也不能把图像孔洞保护直接推成latent干预后的像素不变。

\textbf{输出量化。}按原逐帧最小值阈值决定是否从$[-1,1]$映射，乘255、截断范围、转uint8，采用截断而非四舍五入。保存raw与uint8并核全24帧逐字节对应，再读取固定参考评分。准备期核验发现浮点阈值和负零的库语义差别，已按原计算分别处理；没有看真实分数后改变规则。

\clearpage
\section*{4. 图像混合：完整数据和分母}
以下为图像混合的全部12条新记录。四目标没有筛选；支持区与孔洞区都是旧固定mask，不能解释为真正可见或被遮挡的真值标签。
\begin{center}\small\begin{tabular}{lrrrr}\toprule
强度/目标 & 全图SSE & 全图MSE & 支持区MSE & 孔洞区MSE\\\midrule
0.5/20 & 2,363,633,235 & 0.0365202301 & 0.0365612413 & 0.0358591505\\
0.5/21 & 4,168,341,632 & 0.0644045757 & 0.0561452169 & 0.1254153457\\
0.5/22 & 5,698,872,312 & 0.0880526324 & 0.0773442633 & 0.1326800656\\
0.5/23 & 4,605,301,447 & 0.0711559925 & 0.0644535397 & 0.0970679097\\
0.75/20 & 2,013,502,155 & 0.0311103944 & 0.0308157975 & 0.0358591505\\
0.75/21 & 3,097,530,317 & 0.0478595910 & 0.0373604479 & 0.1254153457\\
0.75/22 & 4,752,138,504 & 0.0734247552 & 0.0592064217 & 0.1326800656\\
0.75/23 & 3,951,095,240 & 0.0610479263 & 0.0517308913 & 0.0970679097\\
1.0/20 & 2,473,419,375 & 0.0382165234 & 0.0383627669 & 0.0358591505\\
1.0/21 & 3,133,839,226 & 0.0484205959 & 0.0379973991 & 0.1254153457\\
1.0/22 & 4,834,069,604 & 0.0746906634 & 0.0607760850 & 0.1326800656\\
1.0/23 & 4,071,700,261 & 0.0629113808 & 0.0540763525 & 0.0970679097\\
\bottomrule\end{tabular}\end{center}
\begin{center}\small\begin{tabular}{lrrr}\toprule
强度 & 全图四帧均值 & 支持等帧/合并均值 & 孔洞等帧/合并均值\\\midrule
0.5 & 0.0650333577 & 0.0586260653 / 0.0576810259 & 0.0977556179 / 0.1086796349\\
0.75 & 0.0533606667 & 0.0447783896 / 0.0440420391 & 0.0977556179 / 0.1086796349\\
1.0 & 0.0560597909 & 0.0478031509 / 0.0471958379 & 0.0977556179 / 0.1086796349\\
\bottomrule\end{tabular}\end{center}
每帧全图331,776像素、995,328个RGB颜色标量；每策略四帧共3,981,312个标量。四目标的支持像素依次312,396、292,217、267,572、263,594；孔洞依次19,380、39,559、64,204、68,182。

区域“等帧”是先算每帧区域误差，再四帧各占四分之一；“合并”按区域实际颜色标量数加权。因为各帧支持面积不同，两种结果可不同，不能互换。空区域用NA，缺失帧不能填零或缩减平均。

全部原始整数SSE、通道分母、空区标记及有符号对比均附CSV/JSON。表中只压缩显示小数位，机器记录保留完整值。旧16行另附，不混进24新行。

\clearpage
\section*{5. 末步融合：完整数据和分母}
以下为末步融合的全部12条新记录。四目标没有筛选；支持区与孔洞区都是旧固定mask，不能解释为真正可见或被遮挡的真值标签。
\begin{center}\small\begin{tabular}{lrrrr}\toprule
强度/目标 & 全图SSE & 全图MSE & 支持区MSE & 孔洞区MSE\\\midrule
0.5/20 & 2,377,007,805 & 0.0367268791 & 0.0377677833 & 0.0199480190\\
0.5/21 & 4,381,005,693 & 0.0676904241 & 0.0619097336 & 0.1103916053\\
0.5/22 & 6,091,937,062 & 0.0941258314 & 0.0836814988 & 0.1376528854\\
0.5/23 & 4,639,729,766 & 0.0716879405 & 0.0670684121 & 0.0895471991\\
0.75/20 & 1,757,491,847 & 0.0271548080 & 0.0279428733 & 0.0144515860\\
0.75/21 & 2,931,721,564 & 0.0452976987 & 0.0369599481 & 0.1068875389\\
0.75/22 & 4,843,245,116 & 0.0748324332 & 0.0578102208 & 0.1457729887\\
0.75/23 & 3,714,051,273 & 0.0573853867 & 0.0492335342 & 0.0889007338\\
1.0/20 & 2,195,754,943 & 0.0339263616 & 0.0352857315 & 0.0120139924\\
1.0/21 & 2,843,983,075 & 0.0439420612 & 0.0365227926 & 0.0987471986\\
1.0/22 & 4,829,395,115 & 0.0746184384 & 0.0587009457 & 0.1409550428\\
1.0/23 & 3,730,615,732 & 0.0576413223 & 0.0493135051 & 0.0898369553\\
\bottomrule\end{tabular}\end{center}
\begin{center}\small\begin{tabular}{lrrr}\toprule
强度 & 全图四帧均值 & 支持等帧/合并均值 & 孔洞等帧/合并均值\\\midrule
0.5 & 0.0675577688 & 0.0626068570 / 0.0615958193 & 0.0893849272 / 0.1029502005\\
0.75 & 0.0511675817 & 0.0429866441 / 0.0422402841 & 0.0890032119 / 0.1041634637\\
1.0 & 0.0525320459 & 0.0449557437 / 0.0443758652 & 0.0853882973 / 0.1009502807\\
\bottomrule\end{tabular}\end{center}
每帧全图331,776像素、995,328个RGB颜色标量；每策略四帧共3,981,312个标量。四目标的支持像素依次312,396、292,217、267,572、263,594；孔洞依次19,380、39,559、64,204、68,182。

区域“等帧”是先算每帧区域误差，再四帧各占四分之一；“合并”按区域实际颜色标量数加权。因为各帧支持面积不同，两种结果可不同，不能互换。空区域用NA，缺失帧不能填零或缩减平均。

全部原始整数SSE、通道分母、空区标记及有符号对比均附CSV/JSON。表中只压缩显示小数位，机器记录保留完整值。旧16行另附，不混进24新行。

\clearpage
\section*{6. 目标20：全部新策略与固定参照}

上排从左到右：实拍参考、旧G0、新图像混合0.5/0.75/1；下排：固定历史投影、旧Gguide、新末步融合0.5/0.75/1。所有小图保持同尺度和完整画面，灰格仅标投影无支持。
\begin{center}\includegraphics[width=\linewidth]{images/target_20_comparison.png}\end{center}
这是为讲解缩小的全图对照，原目录保留24张576像素新图，以及旧批次完整参照。截图或缩略版可帮助定位，细节判断须回到原尺寸。

\textbf{实际人工观察。}图像混合0.5/0.75的显示器、杯子、电话和桌沿有双层轮廓；末步0.5明显软化，0.75仍有杯子/显示器边缘叠加。强度1更接近投影主轮廓，但保留点状破碎纹理、桌沿噪点和边缘残影，不能称正确重建。

观察发生在知道批次/部分数值之后，是团队内非盲描述，不是独立盲评或定量感知分数。可见清晰不保证位置正确；某一张漂亮也不能覆盖同策略其他目标的缺陷。

\clearpage
\section*{7. 目标21：全部新策略与固定参照}

上排从左到右：实拍参考、旧G0、新图像混合0.5/0.75/1；下排：固定历史投影、旧Gguide、新末步融合0.5/0.75/1。所有小图保持同尺度和完整画面，灰格仅标投影无支持。
\begin{center}\includegraphics[width=\linewidth]{images/target_21_comparison.png}\end{center}
这是为讲解缩小的全图对照，原目录保留24张576像素新图，以及旧批次完整参照。截图或缩略版可帮助定位，细节判断须回到原尺寸。

\textbf{实际人工观察。}图像0.5/0.75中显示器和桌面上下错开形成明显重叠；末步0.75仍能看到桌沿双层、下方胶带/书本残影与前景涂抹。强度1的主要物体轮廓较明确，但桌腿/地面区存在破碎或模糊，图像孔洞沿用G0造成接缝。

观察发生在知道批次/部分数值之后，是团队内非盲描述，不是独立盲评或定量感知分数。可见清晰不保证位置正确；某一张漂亮也不能覆盖同策略其他目标的缺陷。

\clearpage
\section*{8. 目标22：全部新策略与固定参照}

上排从左到右：实拍参考、旧G0、新图像混合0.5/0.75/1；下排：固定历史投影、旧Gguide、新末步融合0.5/0.75/1。所有小图保持同尺度和完整画面，灰格仅标投影无支持。
\begin{center}\includegraphics[width=\linewidth]{images/target_22_comparison.png}\end{center}
这是为讲解缩小的全图对照，原目录保留24张576像素新图，以及旧批次完整参照。截图或缩略版可帮助定位，细节判断须回到原尺寸。

\textbf{实际人工观察。}0.5/0.75两族都出现上下两套桌面/显示器轮廓，末步0.75前景仍模糊且残留大尺度叠加。其本帧MSE比旧Gguide更高，不能被四帧均值掩盖。强度1减少部分半透明重影，但仍有投影斑点、断裂桌沿和下部孔洞边界。

观察发生在知道批次/部分数值之后，是团队内非盲描述，不是独立盲评或定量感知分数。可见清晰不保证位置正确；某一张漂亮也不能覆盖同策略其他目标的缺陷。

\clearpage
\section*{9. 目标23：全部新策略与固定参照}

上排从左到右：实拍参考、旧G0、新图像混合0.5/0.75/1；下排：固定历史投影、旧Gguide、新末步融合0.5/0.75/1。所有小图保持同尺度和完整画面，灰格仅标投影无支持。
\begin{center}\includegraphics[width=\linewidth]{images/target_23_comparison.png}\end{center}
这是为讲解缩小的全图对照，原目录保留24张576像素新图，以及旧批次完整参照。截图或缩略版可帮助定位，细节判断须回到原尺寸。

\textbf{实际人工观察。}0.5/0.75有上方投影桌面与下方原生成显示器同时存在的明显重影；末步0.75的下方显示器、鼠标和桌面残影仍在。强度1上方轮廓较鲜明，底部仍有来源于原场景的显示器边缘或模糊残留；地球仪和桌沿点状纹理保留。

观察发生在知道批次/部分数值之后，是团队内非盲描述，不是独立盲评或定量感知分数。可见清晰不保证位置正确；某一张漂亮也不能覆盖同策略其他目标的缺陷。

\clearpage
\section*{10. 为什么MSE低，也可能有两个影子？}

\boxnote{下面全是人工教学数值，不是S86/S87实际实验数据。其作用是解释一种可能，不是给当前重影作唯一归因。}
假设一幅图只有两个位置，亮物体只能在左或右：$A=(1,0)$，$B=(0,1)$。真值各以一半概率出现，已有观测无法区分两者。

直接平均得到$C=(0.5,0.5)$，像两个半亮影子。如果真实是A：
\[\operatorname{MSE}(C,A)=\frac{(0.5-1)^2+(0.5-0)^2}{2}=0.25.\]
真实是B也得到0.25，所以平均风险为0.25。

如果独立抛硬币完整选A或B，选对时误差0，选错时误差$[(1-0)^2+(0-1)^2]/2=1$。各一半，期望误差为0.5。它每次保留一个亮位置，却在这个指标上更差。若有额外信息保证选对当然可得0，但本例没有，不能偷用答案。

逐位置硬拼也未必安全：左边选A、右边选B得到$(1,1)$；反过来得到$(0,0)$。因此“软换硬”本身既不保证正确形状，也不构成创新。

对同一线性空间、量化之前的固定候选$a,b$，令$N$为颜色标量数，$L(x,r)=\|x-r\|^2/N$，$0\leq\alpha\leq1$；随机二选一时以概率$\alpha$选$b$。存在平方展开恒等式：
\[L((1-\alpha)a+\alpha b,r)=(1-\alpha)L(a,r)+\alpha L(b,r)
-\alpha(1-\alpha)\|a-b\|^2/N.\]
最后一项说明平均可降低相对随机二选一的期望平方误差，不要求两候选位置一致；但它不保证胜过更好的端点。不能把该式直接套进50步反馈、非线性VAE与uint8量化的主分数。

Blau与Michaeli的CVPR2018工作讨论感知与失真的关系，本次阅读的是作者v4扩展版。Freirich等NeurIPS2021工作又表明，在正确条件均值、适当端点和最优耦合等前提下，线性混合可能达到最佳折中。因此不能简单说所有混合都错，当前也没有证明这些条件成立。

原文：\url{https://arxiv.org/pdf/1711.06077v4}\par
\url{https://proceedings.neurips.cc/paper_files/paper/2021/file/d77e68596c15c53c2a33ad143739902d-Paper.pdf}

\clearpage
\section*{11. 从结果到创新：该停止什么，保留什么？}
已有普通末端策略达到或低于旧多步引导的本例MSE；“取得该分数必须多步引导”的说法被反例否定。

Gpaste族内三策略的最低MSE为0.0533606667，并列达到者：Gpaste\_l075。

Gterminal族内三策略的最低MSE为0.0511675817，并列达到者：Gterminal\_l075。

all\_six\_new合并六策略的最低MSE为0.0511675817，并列达到者：Gterminal\_l075。

这些是看过参考后的固定有限集合最低值。旧0.25不加入本轮包络，G0原始链仅列一次；没有按目标分别挑强度，也不继续插点寻找更低数字。

\textbf{能保留的贡献是什么？}完成一个可追踪的普通对照，检验了一个会导致过度解释的主张；这属于基线和失败分析证据，不是新生成方法。已有NVS-Solver等已经讨论软融合模糊及排序选择，普通投影、软融合、硬替换或增加强度不能单独包装成新颖性。

\textbf{下一步的选择原则。}停止把最低MSE当唯一目标。只有当更便宜的普通选择对照能回答明确未解的形状问题时才执行；它必须另立协议，不能偷偷加入S87。真正的方法主线需要未用于选规则的场景、独立的请求相机/对象位置参考，检验清晰度与位置正确性是否同时满足。

\textbf{保持未完成标签。}长期、动态、跨场景泛化和新机制验证仍未完成。\texttt{NO\_METHOD\_SELECTED / novelty\_authorization=NONE / new\_method\_validated=false}。这不会抹去普通基线的真实工作，只防止把阶段性证据说成最终论文结论。

\clearpage
\section*{12. 真实成本、复核边界和汇报时的回答}
本轮科学进程总计53.741317秒；VAE构造与权重反序列化计时为0.214885秒（不含此前权重字节读取核验，总进程时间包含这些动作）。3次完整8槽解码，共24个chunk，3次RGB混合；0去噪网络、0新几何、0新编码。旧完整G0链成本不能因为复用而当作不存在，所以不能把本轮派生时长当完整系统端到端速度。
\begin{center}\small\begin{tabular}{lrr}\toprule
策略 & 实际秒数 & 解码chunk数\\\midrule
图像混合 0.5 & 0.031862 & 0\\
末步融合 0.5 & 16.612894 & 8\\
图像混合 0.75 & 0.031210 & 0\\
末步融合 0.75 & 16.544461 & 8\\
图像混合 1.0 & 0.032734 & 0\\
末步融合 1.0 & 16.605043 & 8\\
\bottomrule\end{tabular}\end{center}
科学进程峰值RSS为8.468155 GiB；评分耗时0.436164秒。阶段耗时有包含关系，不重复相加；机器运行时间不是学生工时。

\textbf{复核覆盖到哪里？}团队内不同作者从保存输入以NumPy重算三档clean、Euler、RGB混合、保护区域及24帧量化；没有重新运行VAE来声称外部复现。评分另以按差值直方图累加的整数SSE和精确有理数重算，主判断按整数符号而非显示小数。所有比较都基于封存文件身份。

\textbf{准备期问题没有删除。}接口计数嵌套、每策略回执自引用、浮点阈值和负零语义均在真实运行之前定位，修正程序或核验预期；完整实际失败状态以原回执为准，不能将人工测试写成真实模型运行。

\textbf{证据包。}本目录附可编辑LaTeX、完整24条新CSV、历史16条CSV、实际保存数组、逐帧PNG、执行和评分回执、不同作者审查、固定协议、文献边界与手算讲解。模型权重不重复拷贝。原156页报告保留，新连续版附加本增补，旧章节的时间截点仍是历史记录。

\textbf{老师问：现在最有把握的一句话是什么？}可以清楚说出本轮固定比较的结果，同时说明重影和几何正确性仍需分开检验。

\textbf{老师问：算力还缺什么？}本机能够完成有限派生、分析和小规模验证；更广的多场景、多随机种子、长序列和方法训练需要根据模型与预算再提出GPU需求，不能把尚未测试的显存估算当保证。

\textbf{老师问：为什么没有宣称PhD级新方法？}因为目前证据只支持普通基线与机制排除，最近邻新颖性和独立验证尚未完成。目标可以高，结论必须与证据一致。


\end{document}
